\documentclass{article}
\usepackage[utf8]{inputenc}
\usepackage{amsmath}

\begin{titlepage}
\title{Solución Detallada Ejercicio 10.4 - Econometría}
\author{Análisis de Multicolinealidad Perfecta}
\date{}
\end{titlepage}

\begin{document}

\maketitle

\section{Premisa del Problema}
Se establece que existe una relación lineal exacta entre las variables $X_1$, $X_2$ y $X_3$:
\begin{equation}
\lambda_1 X_{1i} + \lambda_2 X_{2i} + \lambda_3 X_{3i} = 0
\end{equation}
donde al menos uno de los coeficientes $\lambda$ es distinto de cero. Esto implica que cualquier variable puede expresarse como una función exacta de las otras dos. Por ejemplo:
\begin{equation}
X_{1i} = -\frac{\lambda_2}{\lambda_1}X_{2i} - \frac{\lambda_3}{\lambda_1}X_{3i}
\end{equation}

\section{Paso 1: Estimación de Coeficientes de Correlación Parcial}

Los coeficientes de correlación parcial ($r_{12.3}$, $r_{13.2}$, $r_{23.1}$) miden la asociación entre dos variables manteniendo la tercera constante.

En un escenario de multicolinealidad perfecta, la relación es tan "rígida" que, una vez que conocemos los valores de dos variables, la tercera queda totalmente determinada sin error aleatorio.

\subsection{Resultado para $r_{12.3}, r_{13.2}$ y $r_{23.1}$}
Debido a que la relación es lineal y exacta, la correlación entre cualquier par de variables, eliminando el efecto de la tercera, será perfecta. Por lo tanto:
\begin{equation}
|r_{12.3}| = 1, \quad |r_{13.2}| = 1, \quad |r_{23.1}| = 1
\end{equation}
El signo depende de los signos de los coeficientes $\lambda$.

\section{Paso 2: Estimación de Coeficientes de Determinación Múltiple ($R^2$)}

El coeficiente de determinación múltiple $R^2_{y.12}$ mide qué proporción de la varianza de la variable dependiente es explicada por las variables independientes.

\subsection{Análisis de $R^2_{1.23}$}
Aquí $X_1$ actúa como variable "dependiente" frente a $X_2$ y $X_3$. Como vimos en la ecuación (2), $X_1$ es una función lineal exacta de $X_2$ y $X_3$.
\begin{itemize}
    \item \textbf{Lógica:} Si la relación es exacta, no hay residuos ($u_i = 0$).
    \item \textbf{Fórmula:} $R^2 = \frac{ESS}{TSS}$. Si el ajuste es perfecto, la suma de cuadrados explicada (ESS) es igual a la suma de cuadrados totales (TSS).
\end{itemize}

Por lo tanto:
\begin{equation}
R^2_{1.23} = 1
\end{equation}

\subsection{Análisis de $R^2_{2.13}$ y $R^2_{3.12}$}
Siguiendo la misma lógica, como podemos despejar $X_2$ en función de $X_1, X_3$, y $X_3$ en función de $X_1, X_2$:
\begin{equation}
R^2_{2.13} = 1 \quad \text{y} \quad R^2_{3.12} = 1
\end{equation}

\section{Paso 3: Grado de Multicolinealidad}
El grado de multicolinealidad en esta situación es \textbf{Perfecto}.

\subsection{Consecuencias en Econometría}
\begin{enumerate}
    \item \textbf{Inestabilidad de los estimadores:} Los estimadores de MCO (Mínimos Cuadrados Ordinarios) son indeterminados (0/0).
    \item \textbf{Varianza:} Las varianzas y errores estándar de los coeficientes tienden a infinito.
    \item \textbf{Interpretación:} Es imposible aislar el efecto individual de una variable sobre la variable explicada, ya que no se puede variar una manteniendo las otras constantes.
\end{enumerate}

\end{document}
